\begin{exercisechunk}
\item \textbf{Learning noiseless parity from linear equations.}
\label[exercise]{ex:l7-parity-identification}
\exerciseprofile{2}{2}{2}{\exproof\quad\excalculation}{%
Reduce noiseless parity recovery to linear equations over $\mathbb F_2$ and
prove a high-probability recovery guarantee.}
Let $S\subseteq[d]$ be unknown, and let
$(X^{(1)},Y^{(1)}),\ldots,(X^{(n)},Y^{(n)})$ be independent with law $P_S$
as defined in \cref{sec:parity-setup}.

You may use the following standard linear-algebra fact without proof. Given
$A\in\mathbb F_2^{n\times d}$ and $b\in\mathbb F_2^n$, there is an algorithm
that determines in a polynomial number of arithmetic operations whether
$Az=b$ has a unique solution $z\in\mathbb F_2^d$ and returns that solution
when it is unique. If $z$ is one solution, then it is unique if and only if
\[
Aa\ne0
\qquad\text{for every }a\in\mathbb F_2^d\setminus\{0\}.
\]
The usual linear-system theorem over $\mathbb R$ remains valid with
$\mathbb F_2$ in place of $\mathbb R$ because both are fields.
For background, see Sections 2--4 of David Vogan's
\href{https://math.mit.edu/~dav/gauss.pdf}{notes on Gaussian elimination}.

Construct an algorithm that uses a polynomial number of arithmetic operations
and whose output $\widehat S$ satisfies
\[
\P(\widehat S\ne S)<2^{d-n}.
\]
Then prove that, for every $\delta\in(0,1)$, the algorithm recovers $S$ with
probability at least $1-\delta$ whenever
\[
n\ge d+\left\lceil\log_2\frac1\delta\right\rceil.
\]
\begin{hints}
Write
\[
X_j^{(i)}=(-1)^{U_{ij}},
\qquad
Y^{(i)}=(-1)^{V_i},
\qquad
s_j=\I\{j\in S\}.
\]
With arithmetic over $\mathbb F_2$, the examples give
$\mathbf v=\mathbf U\mathbf s$. Apply the stated linear-system procedure. To
bound its failure probability, first fix a nonzero
$a\in\mathbb F_2^d$ and calculate $\P(\mathbf Ua=0)$. Then use a union bound
over the $2^d-1$ possible choices of $a$.
\end{hints}
\end{exercisechunk}
